\section{软件无线电在做什么}

\textbf{软件定义无线电（SDR, Software Defined Radio）}把本应由模拟电路固定的混频、滤波、解调等步骤，尽可能放到数字域用软件完成。
硬件负责射频前端与采样；主机或板载处理器对复数基带样本做同步、均衡、译码。

在 RoboMaster 雷达站场景中，SDR 常用于接收赛场规定的干扰波 / 信息波等波形，解出坐标、密钥等字段，再交给地面视觉融合与裁判上报。
\textbf{方向名「雷达」}指高台感知站；本章的射频硬件与导航章中的\textbf{激光雷达（LiDAR）}无关。

\section{从天线到 IQ}

简化接收链：
\begin{enumerate}
    \item 天线感应射频；
    \item 前端放大、滤波、下变频到中频或零中频；
    \item ADC 采样得到同相（I）与正交（Q）两路，合称 \textbf{IQ} 样本；
    \item 软件对 IQ 做数字滤波、相关、同步与帧解析。
\end{enumerate}

IQ 可看成复信号 $x = I + jQ$。采样率决定可观察的带宽量级；过高采样会受总线（如 USB）吞吐限制，需在带宽与丢样之间折中。

常用教学 / 开发板基于 Analog Devices 的 AD936x 系列收发器（如 ADALM-PLUTO 一类平台）。
主机侧常见软件栈：\textbf{libiio}（Industrial I/O）与 Python 封装 \textbf{pyadi-iio}，用于设置本振频率、采样率、增益，并读取 RX 缓冲。

示意（非赛场参数）：
\begin{tpython}
import adi
sdr = adi.ad9361(uri="ip:192.168.2.1")
sdr.sample_rate = 1_000_000
sdr.rx_lo = 2_400_000_000
iq = sdr.rx()  # complex 数组
\end{tpython}

实战中 URI、网段、天线与功率必须按队内部署文档配置；\textbf{任何发射}都须遵守当地频谱法规与比赛规则。

\section{波形、信道与「分阶段切换」}

赛场无线电往往不是单一载波一直不变，而是多种 profile（中心频、带宽、调制）分阶段出现，例如：
\begin{center}
干扰波一级 $\rightarrow$ 干扰波二级 $\rightarrow$ 信息波
\end{center}
同一块单通道 SDR 在任一时刻只盯一个 profile：确认上一级成功后再切换，避免「同时解两路」带来的复杂度与未验证风险。

工程上把每套射频参数做成 \textbf{radio profile}（本振、采样率、滤波系数等），由状态机驱动切换。
状态机输入可来自地面程序下发的阵营、干扰等级等控制字节。

\section{电磁波解析：从 IQ 到有效帧}

「解析」指从 IQ 流中恢复出协议定义的比特与字段。常见阶段：

\begin{enumerate}
    \item \textbf{前端数字处理}：直流消除、低通 / 匹配滤波、必要时抽取降速；
    \item \textbf{同步}：用已知前导或相关峰找到符号/帧起点（硬判决或软相关）；
    \item \textbf{解调与解映射}：按调制方式得到比特或软信息；
    \item \textbf{组帧与校验}：按长度切出候选帧，做 CRC；失败则滑动搜索下一候选；
    \item \textbf{字段解释}：命令字、坐标、密钥等进入上层（TCP / 融合模块）。
\end{enumerate}

\textbf{软判决}保留幅度/相关度量，而不是过早变成 0/1，往往在低信噪比下更稳。
对关键命令可做长度约束、字段边界搜索，减少盲目组合；CRC 失败的帧必须丢弃。

无硬件时，可用录制的 IQ 短样做离线回放，验证解码器回归——这是射频模块可复现测试的重要手段。

\section{与地面视觉的衔接}

解码成功的坐标帧（示意命令字如坐标上报）经本机 TCP 送给视觉进程后，融合层通常按\textbf{新鲜度}取舍：
无线电新数据优先；过期则回退视觉定位或其它兜底。
非坐标信息（血量、弹药、密钥等）可另槽转发，不阻塞坐标最新值。

通信层细节见本卷「多机通信方法」章；场地图定位与检测流水线见雷达方向其它专项与卷三。

\section{安全与合规*}

\begin{itemize}
    \item 开源或实验配置中的发射功率可能极低，\textbf{不等于}可以随意发射；须遵守法规与赛场管理。
    \item 不要在公开仓库提交真实比赛频点、大体积 IQ、含隐私的录包；用脱敏短样与哈希校验复现。
    \item 射频线缆、供电与接地不良会导致锁失、误码，排障时先分「射频前端 / 解码软件 / 上层通信」。
\end{itemize}

\section{练习建议}

\begin{enumerate}
    \item 用框图默写：天线 $\rightarrow$ IQ $\rightarrow$ 同步 $\rightarrow$ CRC $\rightarrow$ 上层消息。
    \item 说明为何单通道方案常用「分阶段切换」而不是同时解两路。
    \item 在无 SDR 条件下，对一段公开 IQ 样例跑通官方/队内离线解码脚本（若有），记录有效帧数。
\end{enumerate}

延伸阅读：Analog Devices Pluto / AD936x 用户文档；GNU Radio 与 libiio 入门；队内雷达开源中的无线电架构与射频安全说明。
